\chapter{Financial Mathematics}

Everything in this chapter is an application of the geometric series. A
capital under compound interest is a geometric sequence; a stream of equal
payments, accumulated or discounted, is a geometric sum; and the central
formula of the previous chapter,
\[
  \sum_{k=0}^{n-1}q^{k}=\frac{q^{n}-1}{q-1}\qquad(q\neq 1),
\]
turns directly into the end value of a savings plan, the present value of
an annuity, the payment on a loan and (via the limit $n\to\infty$)
the value of a perpetuity.

Time is discrete throughout: interest is credited at the end of each
period, and payments arrive at period ends (or period beginnings, for the
annuity-due variant). Continuous compounding $K_{0}\mathrm{e}^{rt}$
appears exactly once, as a limit that bridges to later volumes;
continuous \emph{discounting} of payment streams by integrals is drilled
in \chref{integrals}{7}. All money amounts are in
euros.

\section{Core Rules}

\begin{rulebox}{Simple vs.\ compound interest}
Let $K_{0}$ be the initial capital and $i$ the annual interest rate
(p.a., written as a decimal). Under \emph{simple interest} only $K_{0}$ earns interest:
\[
  K_{n}=K_{0}(1+n\,i) .
\]
Under \emph{compound interest} credited interest itself
earns interest. With the accumulation factor $q=1+i$,
\[
  K_{n}=K_{0}\,q^{n}=K_{0}(1+i)^{n} .
\]
The capitals $K_{0},K_{1},K_{2},\dots$ form a geometric sequence with
ratio $q$; simple interest grows linearly, compound interest
exponentially.
\end{rulebox}

\begin{rulebox}{Sub-annual compounding and the effective annual rate}
If a nominal rate $j$ p.a.\ is credited in $m$ equal
instalments of $j/m$ per period ($m=4$: quarterly, $m=12$: monthly), then
after $t$ years
\[
  K=K_{0}\left(1+\frac{j}{m}\right)^{mt},
\]
and the \emph{effective annual rate} is the
annually compounded rate producing the same growth:
\[
  i_{\mathrm{eff}}=\left(1+\frac{j}{m}\right)^{m}-1 .
\]
For $m>1$ compounding within the year makes
$i_{\mathrm{eff}}>j$. Competing offers are compared by their effective
rates, never by their nominal rates.
\end{rulebox}

\begin{rulebox}{Continuous compounding (bridge)}
Letting the number of instalments grow without bound and using the limit
$\left(1+\frac{r}{m}\right)^{m}\to\mathrm{e}^{r}$ from the chapter on
sequences,
\[
  K(t)=\lim_{m\to\infty}K_{0}\left(1+\frac{r}{m}\right)^{mt}
      =K_{0}\,\mathrm{e}^{rt} .
\]
This is the largest end value a nominal rate $r$ can generate, and the
form in which interest enters most theoretical models. In this chapter it
serves only as this limit; continuous discounting of entire payment
streams (an integral, not a sum) is drilled in \chref{integrals}{7}.
\end{rulebox}

\begin{rulebox}{Discounting and present value}
Solving $K_{n}=K_{0}q^{n}$ for $K_{0}$ gives the \emph{present value} of a payment $K_{n}$ due in $n$ years:
\[
  K_{0}=\frac{K_{n}}{q^{n}}=K_{n}(1+i)^{-n} .
\]
The factor $q^{-1}=\frac{1}{1+i}$ is the discount factor. Payments due at different dates are made
comparable by discounting all of them to today, this single idea
underlies annuities, loans and project appraisal below.
\end{rulebox}

\begin{rulebox}{Annuity: end value and present value}
An ordinary annuity pays the constant amount
$R$ at the \emph{end} of each of the years $1,\dots,n$. Immediately after
the last payment, payment $k$ has earned interest for $n-k$ years, so the
end value is a geometric sum:
\[
  E_{n}=R\sum_{k=1}^{n}q^{\,n-k}=R\sum_{j=0}^{n-1}q^{\,j}
       =R\,\frac{q^{n}-1}{q-1}=R\,\frac{q^{n}-1}{i} .
\]
Discounting $E_{n}$ back $n$ years gives the present value
\[
  B_{n}=E_{n}\,q^{-n}=R\,\frac{1-q^{-n}}{i} .
\]
For an annuity-due, payments at the
\emph{beginnings} of the years, every payment earns one extra year of
interest: multiply $E_{n}$ and $B_{n}$ by $q$.
\end{rulebox}

\begin{rulebox}{Perpetuity}
Let $n\to\infty$ in the present value $B_{n}=R\,\frac{1-q^{-n}}{i}$.
Since $q>1$ we have $q^{-n}\to 0$ (geometric decay, as in the previous
chapter), so a perpetual annual payment $R$ is worth
\[
  B_{\infty}=\frac{R}{i}
\]
today. Read in both directions: a capital $C$ sustains the perpetual
payment $R=C\,i$, and a perpetual payment $R$ bought for the price $P$
implies the rate $i=R/P$.
\end{rulebox}

\begin{rulebox}{Loan amortization}
A debt $S$ at rate $i$ is repaid by $n$ equal annual payments $A$ if the present value of the payments equals
the debt: $S=A\,\frac{1-q^{-n}}{i}$, hence
\[
  A=S\,\frac{i}{1-q^{-n}}=S\,\frac{i\,q^{n}}{q^{n}-1} .
\]
Each payment splits into interest, $i$ times the debt still outstanding,
and principal, which reduces the debt; as the debt falls,
the interest share shrinks and the principal share grows.
Example: $S=1000$, $i=10\%$, $n=2$:
$A=1000\cdot\frac{0.1\cdot 1.21}{0.21}=576.19$, and the schedule closes
exactly:

\begin{center}
\begin{tabular}{c r r r r}
\hline
Year & Debt at start & Interest & Principal & Debt at end\\
\hline
$1$ & $1000.00$ & $100.00$ & $476.19$ & $523.81$\\
$2$ & $523.81$ & $52.38$ & $523.81$ & $0.00$\\
\hline
\end{tabular}
\end{center}
\end{rulebox}

\begin{rulebox}{Net present value}
A project paying the cash flows $C_{0},C_{1},\dots,C_{n}$ at the ends of
years $0,1,\dots,n$ (typically $C_{0}<0$, today's outlay) has
\emph{net present value}
\[
  \mathrm{NPV}=\sum_{t=0}^{n}\frac{C_{t}}{q^{t}}
\]
at the discount rate $i$ (the opportunity cost of capital). Decision
rule: invest if $\mathrm{NPV}>0$; among mutually exclusive projects
evaluated at the same rate, choose the largest NPV. Undiscounted totals
are meaningless: a euro in year~2 is not a euro today.
\end{rulebox}

\section{Worked Example}

\begin{worked}{A $200{,}000$ EUR mortgage at $4\%$ nominal, compounded monthly, over $20$ years: find the monthly payment and the effective annual rate; then appraise a project paying $400$ for three years against an outlay of $1000$ at $10\%$}
\wstep{Step 1. Convert the quoted rate to a \emph{per-period} rate.}
``$4\%$ nominal, compounded monthly'' means $4\%$ divided into twelve equal
pieces:
\[
  i=\frac{0.04}{12}=0.00\overline{3}\ \text{per month},
  \qquad
  n=20\times12=240\ \text{periods}.
\]
\wtrap{The rate and the number of periods must refer to the \emph{same}
unit of time. Using $i=0.04$ with $n=240$, or $i=0.00\overline{3}$ with
$n=20$, are the two ways this goes wrong, and both produce answers that look
plausible. Write the units next to both numbers.}

\wstep{Step 2. Recognise the cash-flow shape: an ordinary annuity.}
The bank lends $P$ today against $240$ equal payments $R$ at the \emph{end}
of each month. Present value of the payments $=$ amount lent:
\[
  P=R\cdot\frac{1-(1+i)^{-n}}{i}
  \qquad\Longrightarrow\qquad
  R=P\cdot\frac{i}{1-(1+i)^{-n}} .
\]

\wstep{Step 3. Evaluate the discount factor, carefully.}
\[
  (1+i)^{-240}=(1.00\overline{3})^{-240}=0.449927 ,
  \qquad
  1-(1+i)^{-n}=0.550073 .
\]
\[
  R=200\,000\cdot\frac{0.00\overline{3}}{0.550073}
   =\frac{666.667}{0.550073}
   =\boxed{1{,}211.96\ \text{EUR per month}}
\]
\wtrap{Keep six significant figures until the last line. Rounding
$(1+i)^{-240}$ to $0.45$ shifts the payment by about a euro a month --- 
which over $240$ months is a few hundred euros, and in an examination is a
wrong answer.}

\wstep{Step 4. Check by running the formula backwards.}
$1211.96\times\frac{1-(1.00\overline{3})^{-240}}{0.00\overline{3}}
=199\,999.9$, which is $200{,}000$ EUR up to the rounding of the payment.
\checkmark\ Total paid: $240\times1211.96=290{,}871$ EUR, so
$90{,}871$ EUR is interest --- a useful reality check on the size of the
answer.

\wstep{Step 5. The effective annual rate: compounding is not addition.}
Twelve months of $0.00\overline{3}$ do not give $4\%$, because interest
earns interest:
\[
  \mathrm{EAR}=\Bigl(1+\frac{0.04}{12}\Bigr)^{12}-1=1.040742-1=\boxed{4.0742\%}
\]
The gap of $7.4$ basis points is the compounding effect. Continuous
compounding is the limiting case, $\mathrm{e}^{0.04}-1=4.0811\%$, and it is
an upper bound: no compounding frequency can exceed it.
\wtrap{Compare loans by EAR, never by nominal rate. And do not subtract the
$1$ too early --- the ``$-1$'' converts a growth \emph{factor} into a
\emph{rate}, and omitting it turns $4\%$ into $104\%$.}

\wstep{Step 6. The project: discount every flow to today.}
Outlay $1000$ now, then $400$ at the end of each of three years, at $r=10\%$:
\[
  \mathrm{NPV}=-1000+\frac{400}{1.1}+\frac{400}{1.1^{2}}+\frac{400}{1.1^{3}}
  =-1000+363.64+330.58+300.53=\boxed{-5.26\ \text{EUR}}
\]
Negative, so \textbf{reject}. Note how close to zero it is: the project
almost pays for itself, which is why inspection alone cannot decide here
and the calculation is needed.
\wtrap{The first flow is at $t=0$ and is \emph{not} discounted --- exponent
$0$, not $1$. Discounting the outlay as well is the most common error in
NPV, and it flatters every project.}

\wstep{Step 7. Cross-check with the internal rate of return.}
The IRR solves $\mathrm{NPV}(r)=0$; numerically $r^{*}\approx9.70\%$. Since
$9.70\%<10\%$, the project earns less than the required return --- the same
verdict as the NPV, reached independently.
\[
  \mathrm{NPV}>0 \iff \mathrm{IRR}>r
  \quad\text{(for a conventional cash flow: one sign change).}
\]
\wtrap{That equivalence needs the ``one sign change'' proviso. A cash flow
that turns negative again later can have several IRRs or none, and then only
the NPV is meaningful. When the two criteria disagree, rely on the NPV.}

\wstep{Step 8. The three formulas this example used, in one place.}
\[
  \text{annuity PV}=R\,\frac{1-(1+i)^{-n}}{i},
  \qquad
  \text{annuity FV}=R\,\frac{(1+i)^{n}-1}{i},
  \qquad
  \text{perpetuity}=\frac{R}{i} .
\]
All three are geometric series from \chref{foundations}{2}: the perpetuity is the
infinite case, and letting $n\to\infty$ in the first indeed gives $R/i$.
\end{worked}

\section{Practice Problems}

\begin{problems}
  \item A capital of $K_{0}=1000$ is invested for $10$ years at $5\%$
        p.a. Compute the end value under (a) simple interest,
        (b) annual compound interest.
  \item Compute the end value of $K_{0}=2500$ after $5$ years at $4\%$
        p.a.\ compound interest.
  \item What is the present value of a payment of $10\,000$ due in
        $8$ years, at $5\%$ p.a.?
  \item A zero-coupon bond repays $5000$ in $6$ years. At a market rate
        of $2\%$ p.a., what price (present value) is fair today?
  \item A bank offers a nominal rate of $8\%$ p.a.\ compounded
        quarterly. Compute the effective annual rate and decide whether
        a saver should prefer this account or one paying $8.25\%$ with
        annual compounding.
  \item Compute the effective annual rate of a nominal $12\%$ p.a.\
        compounded monthly.
  \item How long does it take to double a capital at $5\%$ p.a.\
        compound interest? (Solve $1.05^{n}=2$ with logarithms.)
  \item How long does it take to double a capital at $2\%$ p.a.\
        compound interest?
  \item $K_{0}=1000$ is invested for $10$ years at $r=5\%$ p.a.\ with
        \emph{continuous} compounding, $K(t)=K_{0}\mathrm{e}^{rt}$.
        Compute the end value and compare with annual compounding from
        Problem~1.
  \item Savings plan: you save $200$ per month; for clean numbers assume
        the whole $R=2400$ is deposited at the end of each year. At
        $5\%$ p.a., what has the plan grown to immediately after the
        $10$th deposit?
  \item Repeat Problem~10 with the deposits made at the
        \emph{beginning} of each year (annuity-due).
  \item An ordinary annuity pays $R=1000$ at the end of each year for
        $5$ years; $i=10\%$. Compute its present value.
  \item Perpetuity: what is the present value of $R=500$ per year
        forever at $i=4\%$?
  \item A perpetual bond paying $1500$ per year costs $25\,000$. Which
        interest rate does this price imply?
  \item A debt of $S=100\,000$ is to be repaid by equal annual payments
        over $20$ years at $5\%$ p.a. Compute the annuity payment.
  \item A loan of $10\,000$ at $10\%$ p.a.\ is repaid by three equal
        annual payments. Compute the payment and set up the full
        amortization schedule (interest, principal and remaining debt
        for each of the three years).
  \item Two projects each cost $1000$ today. Project~A returns $600$ at
        the end of each of years $1$ and $2$; project~B returns a single
        $1250$ at the end of year~$2$. Compare the net present values at
        $i=10\%$: which project should be chosen?
  \item A project costs $2000$ today and returns $800$ at the end of
        each of years $1$--$3$. Compute the NPV at $i=5\%$. Is the
        project worthwhile?
  \item Three deposits of $100$, made at the ends of years $1,2,3$, have
        grown to exactly $331$ immediately after the last deposit. Use
        the geometric sum to set up an equation for $q=1+i$ and solve
        for the interest rate $i$.
  \item $K_{0}=1000$ is invested for $3$ years at $4\%$ p.a.\ and then
        for $2$ further years at $10\%$ p.a.\ (compound interest
        throughout). Compute the end value.
\end{problems}

\section{Solutions}

Money is rounded to two decimal places at the final step (and, in
amortization schedules, to the cent in every row); growth and discount
factors are carried at full precision.

\begin{solutions}
  \item (a) Simple interest: $K_{n}=K_{0}(1+n\,i)
        =1000\,(1+10\cdot 0.05)=1000\cdot 1.5=1500.00$.\\
        (b) Compound interest: $K_{n}=K_{0}(1+i)^{n}
        =1000\cdot 1.05^{10}=1000\cdot 1.6288946=1628.89$.
        Compounding earns $128.89$ more.

  \item $K_{n}=K_{0}(1+i)^{n}=2500\cdot 1.04^{5}$. Since
        $1.04^{5}=1.2166529$,
        $K_{5}=2500\cdot 1.2166529=3041.63$.

  \item $K_{0}=K_{n}(1+i)^{-n}=\dfrac{10\,000}{1.05^{8}}$. Since
        $1.05^{8}=1.4774554$,
        $K_{0}=\dfrac{10\,000}{1.4774554}=6768.39$.

  \item $K_{0}=\dfrac{5000}{1.02^{6}}$. Since $1.02^{6}=1.1261624$,
        $K_{0}=\dfrac{5000}{1.1261624}=4439.86$.

  \item $i_{\mathrm{eff}}=\left(1+\frac{j}{m}\right)^{m}-1
        =\left(1+\frac{0.08}{4}\right)^{4}-1=1.02^{4}-1$. Since
        $1.02^{4}=1.08243216$,
        $i_{\mathrm{eff}}=0.08243216\approx 8.2432\%$. As
        $8.2432\%<8.25\%$, the account with annual compounding is
        (slightly) better for the saver.

  \item $i_{\mathrm{eff}}=\left(1+\frac{0.12}{12}\right)^{12}-1
        =1.01^{12}-1=1.1268250-1=0.1268250\approx 12.6825\%$.

  \item $1.05^{n}=2\Rightarrow n\ln 1.05=\ln 2\Rightarrow
        n=\dfrac{\ln 2}{\ln 1.05}=\dfrac{0.6931472}{0.0487902}
        \approx 14.21$ years. Since interest is credited only at year
        ends, the capital has actually doubled after $15$ full years
        ($1.05^{14}=1.9799<2<2.0789=1.05^{15}$).

  \item $n=\dfrac{\ln 2}{\ln 1.02}=\dfrac{0.6931472}{0.0198026}
        \approx 35.00$ years. (Because $n=35.0028>35$, with annual
        crediting the capital first exceeds the double after $36$ full
        years: $1.02^{35}\approx 1.9999<2$.)

  \item $K(10)=K_{0}\mathrm{e}^{rt}=1000\,\mathrm{e}^{0.05\cdot 10}
        =1000\,\mathrm{e}^{0.5}=1000\cdot 1.6487213=1648.72$.
        Continuous compounding beats the annual value $1628.89$ from
        Problem~1 by $19.83$.

  \item Monthly saving $200$ means $R=12\cdot 200=2400$ per year. End
        value of the ordinary annuity:
        \[
          E_{10}=R\,\frac{q^{10}-1}{i}
                =2400\cdot\frac{1.05^{10}-1}{0.05}
                =2400\cdot\frac{0.6288946}{0.05}
                =2400\cdot 12.5778925=30\,186.94 .
        \]

  \item Annuity-due: every deposit earns one extra year, so
        $E_{10}^{\mathrm{due}}=q\,E_{10}
        =1.05\cdot 30\,186.94=31\,696.29$.

  \item $B_{n}=R\,\dfrac{1-q^{-n}}{i}
        =1000\cdot\dfrac{1-1.1^{-5}}{0.1}$. Since $1.1^{5}=1.61051$,
        $1.1^{-5}=0.6209213$ and
        \[
          B_{5}=1000\cdot\frac{1-0.6209213}{0.1}
               =1000\cdot 3.7907868=3790.79 .
        \]

  \item $B_{\infty}=\dfrac{R}{i}=\dfrac{500}{0.04}=12\,500.00$.

  \item From $B_{\infty}=R/i$: $i=\dfrac{R}{B_{\infty}}
        =\dfrac{1500}{25\,000}=0.06=6\%$.

  \item $A=S\,\dfrac{i\,q^{n}}{q^{n}-1}
        =100\,000\cdot\dfrac{0.05\cdot 1.05^{20}}{1.05^{20}-1}$. Since
        $1.05^{20}=2.6532977$,
        \[
          A=100\,000\cdot\frac{0.05\cdot 2.6532977}{1.6532977}
           =100\,000\cdot\frac{0.1326649}{1.6532977}
           =100\,000\cdot 0.0802426=8024.26 .
        \]

  \item Payment: $A=S\,\dfrac{i\,q^{n}}{q^{n}-1}
        =10\,000\cdot\dfrac{0.1\cdot 1.331}{0.331}
        =\dfrac{1331}{0.331}=4021.15$. Schedule (interest $=10\%$ of the
        debt at start, principal $=A-{}$interest, all rounded to the
        cent):

        \begin{center}
        \begin{tabular}{c r r r r}
        \hline
        Year & Debt at start & Interest & Principal & Debt at end\\
        \hline
        $1$ & $10\,000.00$ & $1000.00$ & $3021.15$ & $6978.85$\\
        $2$ & $6978.85$ & $697.89$ & $3323.26$ & $3655.59$\\
        $3$ & $3655.59$ & $365.56$ & $3655.59$ & $0.00$\\
        \hline
        \end{tabular}
        \end{center}

        The interest share falls ($1000.00$, $697.89$, $365.56$) while
        the principal share grows, and the last payment
        $365.56+3655.59=4021.15$ closes the loan exactly.

  \item $\mathrm{NPV}=\sum_{t}C_{t}q^{-t}$ with $q=1.1$:
        \[
          \mathrm{NPV}_{A}=-1000+\frac{600}{1.1}+\frac{600}{1.1^{2}}
          =-1000+545.45+495.87=41.32 ,
        \]
        \[
          \mathrm{NPV}_{B}=-1000+\frac{1250}{1.1^{2}}
          =-1000+\frac{1250}{1.21}=-1000+1033.06=33.06 .
        \]
        Both are worthwhile, but $\mathrm{NPV}_{A}>\mathrm{NPV}_{B}$:
        choose project~A, even though B's undiscounted total
        ($1250$) exceeds A's ($1200$), A pays earlier.

  \item The three equal returns form an ordinary annuity, so
        \[
          \mathrm{NPV}=-2000+800\cdot\frac{1-1.05^{-3}}{0.05}
          =-2000+800\cdot\frac{1-0.8638376}{0.05}
          =-2000+800\cdot 2.7232480 ,
        \]
        using $1.05^{3}=1.157625$. Hence
        $\mathrm{NPV}=-2000+2178.60=178.60>0$: the project is
        worthwhile.

  \item End value of the three deposits (geometric sum):
        $100\,(q^{2}+q+1)=331$, i.e.\ $q^{2}+q+1=3.31$, so
        $q^{2}+q-2.31=0$. The quadratic formula gives
        \[
          q=\frac{-1+\sqrt{1+4\cdot 2.31}}{2}
           =\frac{-1+\sqrt{10.24}}{2}=\frac{-1+3.2}{2}=1.1
        \]
        (the negative root is meaningless). Hence $i=q-1=0.1=10\%$.
        Check: $100\,(1.21+1.1+1)=331$.

  \item Growth factors multiply:
        $K_{5}=K_{0}\cdot 1.04^{3}\cdot 1.1^{2}
        =1000\cdot 1.124864\cdot 1.21=1361.09$.
\end{solutions}

\section{Additional Problems}

\begin{problems}
  \item Compute the end value of $K_{0}=8000$ after $4$ years at $5\%$
        p.a.\ compound interest.
  \item $K_{0}=5000$ is invested for $10$ years at $6\%$ p.a. Compute
        the end value under simple and under compound interest, and the
        difference.
  \item What is the present value of $20\,000$ due in $10$ years at
        $10\%$ p.a.?
  \item A nominal rate of $4\%$ p.a.\ is compounded quarterly. Compute
        the effective annual rate and compare with an account paying
        $4.05\%$ annually: which should a saver choose?
  \item How long does it take to double a capital (a) at $10\%$ p.a.\
        with annual compounding, (b) at $r=10\%$ with continuous
        compounding $K(t)=K_{0}\mathrm{e}^{rt}$?
  \item A savings plan deposits $R=1000$ at the end of each year for
        $5$ years at $4\%$ p.a. Compute the end value immediately after
        the last deposit.
  \item A pension pays $6000$ at the end of each year for $10$ years;
        $i=5\%$. Compute its present value.
  \item An endowment is to fund a prize of $3000$ per year forever at
        $i=6\%$. How much capital is required today?
  \item A loan of $50\,000$ at $4\%$ p.a.\ is repaid by five equal
        annual payments. Compute the payment.
  \item A project costs $5000$ today and returns $2000$ at the end of
        each of years $1$--$3$. Compute the NPV at $i=10\%$ and decide.
        Compare with the undiscounted total of the returns.
\end{problems}

\section{Additional Solutions}

\begin{solutions}
  \item $K_{n}=K_{0}(1+i)^{n}=8000\cdot 1.05^{4}
        =8000\cdot 1.21550625=9724.05$.

  \item Simple: $K_{n}=K_{0}(1+n\,i)=5000\,(1+10\cdot 0.06)
        =5000\cdot 1.6=8000.00$. Compound:
        $K_{n}=5000\cdot 1.06^{10}=5000\cdot 1.7908477=8954.24$.
        Compounding earns $8954.24-8000.00=954.24$ more.

  \item $K_{0}=\dfrac{20\,000}{1.1^{10}}=\dfrac{20\,000}{2.5937425}
        =7710.87$.

  \item $i_{\mathrm{eff}}=\left(1+\frac{0.04}{4}\right)^{4}-1
        =1.01^{4}-1=1.04060401-1=0.04060401\approx 4.0604\%$. As
        $4.0604\%>4.05\%$, the quarterly account is (slightly) better
        for the saver.

  \item (a) $1.1^{n}=2\Rightarrow n=\dfrac{\ln 2}{\ln 1.1}
        =\dfrac{0.6931472}{0.0953102}\approx 7.27$ years.\\
        (b) $\mathrm{e}^{0.1t}=2\Rightarrow t=\dfrac{\ln 2}{0.1}
        =6.93$ years, continuous compounding doubles faster.

  \item $E_{5}=R\,\dfrac{q^{5}-1}{i}
        =1000\cdot\dfrac{1.04^{5}-1}{0.04}
        =1000\cdot\dfrac{0.2166529}{0.04}
        =1000\cdot 5.4163226=5416.32$.

  \item $B_{10}=R\,\dfrac{1-q^{-10}}{i}
        =6000\cdot\dfrac{1-1.05^{-10}}{0.05}$. Since
        $1.05^{10}=1.6288946$, $1.05^{-10}=0.6139133$ and
        \[
          B_{10}=6000\cdot\frac{1-0.6139133}{0.05}
                =6000\cdot 7.7217349=46\,330.41 .
        \]

  \item $B_{\infty}=\dfrac{R}{i}=\dfrac{3000}{0.06}=50\,000.00$.

  \item $A=S\,\dfrac{i}{1-q^{-n}}
        =50\,000\cdot\dfrac{0.04}{1-1.04^{-5}}$. Since
        $1.04^{5}=1.2166529$, $1.04^{-5}=0.8219271$ and
        \[
          A=50\,000\cdot\frac{0.04}{0.1780729}
           =50\,000\cdot 0.2246271=11\,231.36 .
        \]

  \item $\mathrm{NPV}=-5000+2000\cdot\dfrac{1-1.1^{-3}}{0.1}$. Since
        $1.1^{3}=1.331$, $1.1^{-3}=0.7513148$ and
        \[
          \mathrm{NPV}=-5000+2000\cdot\frac{1-0.7513148}{0.1}
                      =-5000+2000\cdot 2.4868520
        \]
        \[
          \phantom{\mathrm{NPV}}=-5000+4973.70=-26.30<0 :
        \]
        reject the project. The undiscounted total $3\cdot 2000=6000$
        exceeds the outlay $5000$, but discounting at $10\%$ turns the
        surplus into a loss, totals without discounting mislead.
\end{solutions}
